        \paragraph{Why the scalar substrate is bosonic and $\theta=\pi$ is the fermionic $\mathbb{Z}_2$ choice.}
        The two boundary classes are not arbitrary continuous tunings.  In an
        $S^2\simeq\mathbb{C}P^1$ order-parameter theory they can be identified with the two values
        of a discrete topological/statistical sector.

        \textbf{[DERIVED NEGATIVE] Scalar $U(1)$ substrate.}
        A single-component order parameter $\psi=|\psi|e^{i\varphi}$ is single-valued only if
        \begin{equation}
            \oint \mathrm{d}\varphi \in 2\pi\mathbb{Z},
        \end{equation}
        so its scalar phase is periodic:
        \begin{equation}
            e^{i\oint \mathrm{d}\varphi}=+1
        \end{equation}
        on every closed loop.  This condition allows integer windings, but it does not force
        the longitudinal framing class to be odd.  With the positive stiffness
        energy~\eqref{eq:E_chi},
        \begin{equation}
            E_\chi = E_0+\frac{2\pi^2 A_\parallel}{L}\,\chi^2,\qquad A_\parallel>0,
        \end{equation}
        minimization over the ordinary scalar class $\chi\in\mathbb{Z}$ selects
        \begin{equation}
            \chi_{\rm ground}=0 .
        \end{equation}
        Hence a purely scalar SST substrate has no intrinsic mechanism for the odd core-phase
        sector required by fermions.  In this operational sense the scalar substrate is bosonic:
        it admits the trivial periodic carrier and does not generate the $pq\pm1$ lepton ladder.

        \textbf{[DERIVED / ALLOWED] $S^2\simeq\mathbb{C}P^1$ admits a $\mathbb{Z}_2$ statistical sector.}
        The matter target is not a scalar $U(1)$ phase but the director/Hopfion target
        \begin{equation}
            M_{\rm SST}=S^2\simeq\mathbb{C}P^1 .
        \end{equation}
        The relevant homotopy data are
        \begin{equation}
            \pi_3(S^2)=\mathbb{Z},\qquad
            \pi_4(S^2)=\mathbb{Z}_2 .
        \end{equation}
        The first class gives the static Hopf charge,
        \begin{equation}
            Q_H = SL_{\rm phys},
        \end{equation}
        while the second allows a $\mathbb{Z}_2$ topological/statistical phase for spacetime
        loops of soliton configurations~\cite{WilczekZee1983}.  More carefully, the fermionic
        sign is not simply a local scalar factor $e^{i\theta Q_H}$ in the static Hopf charge.
        It is the sign of the relevant $\mathbb{Z}_2$ spacetime class,
        \begin{equation}
            \exp(iS_\theta)=(-1)^\nu,\qquad
            \nu\in\mathbb{Z}_2 ,
        \end{equation}
        whose representative may be read, for the relevant $2\pi$ rotation or exchange process,
        as the parity class of the Hopf sector.  Thus the two possible statistical sectors are
        \begin{align}
            \theta=0 &: \quad \exp(iS_\theta)=+1
                \quad\Rightarrow\quad \text{bosonic / periodic sector},\\
            \theta=\pi &: \quad \exp(iS_\theta)=(-1)^\nu
                \quad\Rightarrow\quad \text{fermionic / spinorial sector for odd Hopf parity}.
        \end{align}
        Identifying the spinorial sector with the longitudinal core-phase boundary class gives
        \begin{equation}
            \theta=\pi
            \quad\Longleftrightarrow\quad
            \chi\in2\mathbb{Z}+1 .
        \end{equation}
        Then positive stiffness no longer minimizes over all integers but over the odd class:
        \begin{equation}
            \chi\in2\mathbb{Z}+1
            \quad\Rightarrow\quad
            \chi_{\rm ground}=\pm1 .
        \end{equation}
        Since
        \begin{equation}
            Q_H=SL_{\rm phys}=SL_{\rm tor}+\chi=pq+\chi,
        \end{equation}
        and $pq=2q$ is even for the lepton family $T(2,q)$, the $\theta=\pi$ sector gives
        \begin{equation}
            SL_{\rm phys}(T(2,q))=2q\pm1,
        \end{equation}
        i.e. the odd Hopf/self-linking carrier required for fermionic leptons.

        \textbf{[POSIT, pinned] $\theta=\pi$ is selected by observation, not derived.}
        The framed-curve geometry and the quadratic stiffness energy do not prefer
        $\theta=\pi$ over $\theta=0$.  The value $\theta=\pi$ is therefore a discrete
        superselection input pinned by the empirical existence of fermions.  It is naturally
        identified with the same $\mathbb{Z}_2$ datum appearing in half-quantum vortices of
        multi-component condensates~\cite{Volovik2003,Autti2016} and in the half/full
        circulation assignment of the magnetic-moment sector (Route~B).  With that
        identification, fermion statistics, the $pq\pm1$ ladder, and the Route~B origin of
        $g=2$ are no longer independent posits but three faces of one $\mathbb{Z}_2$ sector.

        \textbf{[CRITICAL NOTE]}
        The statement ``$\pi_4(S^2)=\mathbb{Z}_2$ guarantees a fermionic Hopfion'' is too strong.
        The homotopy group shows that a $\mathbb{Z}_2$ spacetime class exists, but hopfion
        statistics are set by the presence and value of the corresponding topological/statistical
        term.  In modern language this classification can involve bordism/TQFT data rather than
        only bare homotopy groups~\cite{ChenTanizaki2023}.  What is solid here is the reduction:
        \begin{equation}
            \theta=0 \Rightarrow \text{periodic/bosonic substrate},\qquad
            \theta=\pi \Rightarrow \text{spinorial/fermionic substrate}.
        \end{equation}
        The remaining open problem is to derive why the SST substrate realizes the
        $\theta=\pi$ sector rather than the $\theta=0$ sector.
